% =========================================================
% Proof: Cauchy Sequences Are Bounded
% =========================================================

\newpage
\phantomsection
\label{prf:cauchy-bounded-q}

\begin{remark*}[Return]
\hyperref[thm:cauchy-bounded-q]{Return to Theorem~\ref*{thm:cauchy-bounded-q}.}
\end{remark*}

\begin{theorem*}[Cauchy sequences are bounded]
Every Cauchy sequence of rational numbers is bounded.
\end{theorem*}

\begin{proof}[Proof (Professional Standard)]
Let $(x_n)$ be a Cauchy sequence of rational numbers.
Apply the Cauchy condition with $\varepsilon = 1$: there exists
$N \in \mathbb{N}$ such that $|x_m - x_n| < 1$ for all $m, n \geq N$.
Taking $n = N + 1$, we get $|x_m - x_{N+1}| < 1$ for all $m \geq N$,
so $|x_m| \leq |x_m - x_{N+1}| + |x_{N+1}| < 1 + |x_{N+1}|$
for all $m \geq N$.
Define
\[
\Gamma = \max\{|x_1|, \ldots, |x_N|, 1 + |x_{N+1}|\}.
\]
Then $|x_n| \leq \Gamma$ for all $n \in \mathbb{N}$.
\end{proof}

\begin{proof}[Proof (Detailed Learning)]
\textbf{Step 1: Apply the Cauchy condition with $\varepsilon = 1$.}
By \hyperref[def:cauchy-rational]{Definition~\ref*{def:cauchy-rational}},
there exists $N \in \mathbb{N}$ such that
$|x_m - x_n| < 1$ for all $m, n \geq N$.

\textbf{Step 2: Anchor the tail to a single term.}
Fix $n = N + 1$.  For every $m \geq N$ we have $n = N + 1 \geq N$, so
\[
|x_m - x_{N+1}| < 1.
\]

\textbf{Step 3: Bound each tail term via the triangle inequality.}
For any $m \geq N$:
\[
|x_m| = |x_m - x_{N+1} + x_{N+1}|
\leq |x_m - x_{N+1}| + |x_{N+1}|
< 1 + |x_{N+1}|.
\]

\textbf{Step 4: Handle the initial segment.}
The terms $x_1, \ldots, x_N$ are finitely many, so
$\max\{|x_1|, \ldots, |x_N|\}$ is a well-defined rational number.

\textbf{Step 5: Combine.}
Set
\[
\Gamma = \max\{|x_1|, \ldots, |x_N|, 1 + |x_{N+1}|\}.
\]
For $n \leq N$: $|x_n| \leq \Gamma$.
For $n > N$ (which implies $n \geq N$): $|x_n| < 1 + |x_{N+1}| \leq \Gamma$.
Therefore $|x_n| \leq \Gamma$ for all $n \in \mathbb{N}$.

\textbf{Step 6: Conclude.}
By \hyperref[def:bounded-sequence-q]{Definition~\ref*{def:bounded-sequence-q}},
$(x_n)$ is bounded.
\end{proof}

\begin{remark*}[Proof structure]
The architecture is identical to the proof that convergent sequences are bounded,
with one substitution: instead of anchoring the tail to the limit $a$, we anchor
it to the single term $x_{N+1}$.  The Cauchy condition controls
$|x_m - x_{N+1}|$ for all $m \geq N$ in exactly the same way that convergence
controlled $|x_m - a|$.  This structural parallel is not a coincidence: the
Cauchy condition is convergence without a named limit.
\end{remark*}

\begin{remark*}[Dependencies]
\hyperref[def:cauchy-rational]{Cauchy sequence in $\mathbb{Q}$},
\hyperref[def:bounded-sequence-q]{Bounded sequence},
the triangle inequality in $\mathbb{Q}$.
\end{remark*}
